% 附录 A-H
% 依赖：class 已提供 tabularx、array、amsmath；本文件不新增宏包。
% 表格采用固定标题（表 A1、表 B1 等），与正文表号解耦。
% 正文引用：附录 A-H；表 A1、B1-B3、C1、D1-D5、E1-E2、F1-F3、G1、H1。

\appendix
\ctexset{section={number=附录~\Alph{section}, aftername={\quad}}}

\newcommand{\apptabtitle}[1]{\par\medskip\noindent\textbf{#1}\par\nopagebreak[3]}
\newcommand{\apptabnote}[1]{\par\vspace{2pt}{\footnotesize #1}\par\medskip}
\newcommand{\apptabhead}[1]{\textbf{#1}}

\section{题面要求对照}

本附录逐条对照题面的 12 条要求，并依据仓库内可打开的证据登记状态。第 6、7 条的防护与同规程复测证据见各自条目。第 9 条的可选项分列为优化级别、架构与同机网络三组。

\begin{center}
{\footnotesize
\renewcommand{\arraystretch}{1.05}
\begin{tabularx}{\textwidth}{|c|>{\raggedright\arraybackslash}p{3.45cm}|>{\raggedright\arraybackslash}X|>{\raggedright\arraybackslash}p{1.55cm}|c|>{\raggedright\arraybackslash}p{2.85cm}|}
\hline
\apptabhead{序号} & \apptabhead{要求} & \apptabhead{本作品的做法} & \apptabhead{章节} & \apptabhead{状态} & \apptabhead{证据} \\
\hline
1 & 实现或接入符合规范的 ML-KEM 密钥生成、封装与解封装，用公开或主办方向量验证。不得改安全语义换常数时间。 & 三条实现路径接入统一接口。公开向量、修复等价与隐式拒绝测试。M1 标注为 FIPS 203 初始公开草案（ipd），M2 标注为历史 Kyber768。 & 3.1、3.4、4.2 & 满足 & A-33 至 A-35 \\
\hline
2 & 建立可重复的高精度计时框架，说明处理器、操作系统、编译器、优化级别、频率管理、测试次数与异常样本处理。 & 单调时钟计时。锁频控制器 20 次探针与退出恢复核对。28 项环境表。原始长尾全留，失败样本单列。 & 3.5、3.9、4.1 & 满足 & A-29 至 A-31 \\
\hline
3 & 至少两类有明确区别的解封装输入，重复次数足够。输入设计服务于泄漏分析。 & 主类对为 message-zero 与 message-ones 等长合法密文。辅以位翻转、随机、格式破损与截断。每类 10000 次、三次运行。 & 3.1、4.1、4.8 & 满足 & A-01、A-03、A-28 \\
\hline
4 & 用统计方法判断不同类别的执行时间是否显著可区分。均值、方差、分位数、置信区间与检验并用，不得凭少量样本。 & Welch 检验、Holm 校正、逆方差固定效应合并、I\textsuperscript{2} 与 MDE。RTLF、dudect、Mona 作族外交叉。 & 3.6、4.3、4.7 & 满足 & A-09 至 A-16 \\
\hline
5 & 定位泄漏来源，说明其与解封装流程的对应关系。 & 反汇编、Memcheck、varlat、Microwalk、Binsec。3 实现 × 5 步骤 × 6 类 = 90 格逐格判定，5 格存在、75 格不存在、10 格不适用、0 格未检查。 & 3.7、4.4、4.5 & 满足 & A-17 至 A-23 \\
\hline
6 & 对已发现问题采用常数时间比较、掩码选择、固定循环、统一错误处理或消除敏感分支等方式修改。不得通过日志、异常、返回码泄露私钥、中间值或验证结果。 & 三层防护闭环：M1/M2 接入上游标准补丁并做二进制产物级消融核验；M3/M4 人工夹具自主改造，其中 M4 完成有界形式化验证，M3/M4 均有统计复测；平台 1,048 个输出与返回调用点完成污点审计并修复 1 处退出码侧信道。122 个调用点接触秘密，未发现其他秘密泄露；同规程复测见第 7 条。审计不含 vendor 树。 & 3.8、4.6 & 满足 & A-18、A-23、A-36 \\
\hline

\end{tabularx}
}
\apptabnote{表 A1　题面 12 条对照（一）。状态以仓库内可打开的证据文件为判据，不以文稿措辞替代实验完成。}
\end{center}

\begin{center}
{\footnotesize
\renewcommand{\arraystretch}{1.05}
\begin{tabularx}{\textwidth}{|c|>{\raggedright\arraybackslash}p{3.45cm}|>{\raggedright\arraybackslash}X|>{\raggedright\arraybackslash}p{1.55cm}|c|>{\raggedright\arraybackslash}p{2.85cm}|}
\hline
\apptabhead{序号} & \apptabhead{要求} & \apptabhead{本作品的做法} & \apptabhead{章节} & \apptabhead{状态} & \apptabhead{证据} \\
\hline
7 & 防护后在与防护前相同或等价的测试环境下重采数据，对比不同输入类别的时间分布与统计结果，说明泄露是否明显降低。同时验证功能正确性。 & M1/M2 的上游补丁前后按同一规程复测；M3/M4 人工夹具自主改造前后在定位与计时层作同规程对照。平台调用点审计与退出码修复后复核，公开向量与语义测试通过。 & 4.6、4.9 & 满足 & A-01、A-02、A-06、A-07、A-35 \\
\hline
8 & 给出防护前后解封装平均时间、典型分位数与性能开销。不得牺牲安全性或正确性。 & 25 行按实现版本和操作汇总、11 行成对开销；M1 修复 +1.51\%，M2 修复 $-2.92\%$，M4 比较夹具 $+26.21\%$ 与 $+6.85\%$。 & 4.9 & 满足 & A-25 \\
\hline
9 & 至少在一种软件环境完成完整实验。高水平作品可增加优化级别、CPU 架构或本地与网络场景。 & 12 格优化级别、12 格 aarch64、同机网络 5 档 × 3 种实现 × 2 次。aarch64 为机制对照，网络为同机模拟。 & 4.10、4.11 & 满足 & A-26 至 A-28 \\
\hline
10 & 明确攻击者可控输入、可观测信息、本地权限与攻击条件。区分本地高精度计时与远程网络时间攻击的边界。 & 本地白盒模型、可控输入与观测量见第 1.2 节；适用范围与未验证项见第 6.2 节。 & 1.2、3.2、6.2 & 满足 & A-28、C-05、C-06 \\
\hline
11 & 提供测试程序、数据生成方法、实验参数、编译选项与统计过程，使评测人员可复现。 & 一键入口分 quick 与 full 两档。干净 clone 263 项测试通过。vendor 取件两层 SHA-256。逐次运行台账与原始长表随附。 & 4.13、附录 G & 满足 & A-37 \\
\hline
12 & 至少在自动化检测、敏感路径定位、常数时间实现、统计方法、低开销防护或跨平台验证之一取得改进，并用可复现实验说明效果。 & 分层证据与 MDE。预注册锁频。90 格与 9 × 5 工具矩阵。产物与架构依赖对照。 & 1.4、4.12、5.1、6.1 & 满足 & A-22、A-24、A-26、A-27 \\
\hline
\end{tabularx}
}
\apptabnote{表 A1　题面 12 条对照（二）。状态按仓库内可打开的证据文件判定；第 6、7 条的防护和复测证据均列于同行。}
\end{center}

题面 12 条要求均有对应实现与证据，状态全部为“满足”。防护工作形成三层完整闭环：对 M1、M2 接入标准补丁并核验二进制产物；对 M3、M4 人工夹具实施常数时间改造，以统计复测覆盖两个夹具，并对 M4 比较内核作有界形式化验证；平台对 1,048 个调用点作静态污点审计，修复 1 处退出码侧信道。90 格排查发现的 M2 同源消息编码分支作为方法超额发现记录，不改变第 6、7 条的满足状态。第 9 条的优化级别、架构与同机网络实验均已执行，各自的测量范围按表中说明。

状态判定依据是可打开的原始数据、运行记录与产物摘要。第 9 条的可选项按覆盖范围列示：优化级别有完整 12 格，aarch64 提供机制层 12 格，网络实验覆盖同机模拟 5 档；环境范围见第 1.2 节与第 6.2 节。

12 条分为基础、防护与差异三层：第 1 至 5 条覆盖功能、测量、输入、统计与定位；第 6 至 8 条覆盖防护修改、同规程复测与性能开销；第 9 至 12 条覆盖实验环境、威胁模型、复现与改进。每项结论均回指相应证据，适用范围在对应条目中说明。

证据列用的是台账编号，不是文件路径。A 类编号可在附录 B 至 F 或正文对应节找到原始记录。C 类编号是未执行项，只在边界段引用。编号与章节号并列，是为了让读者既能按章读，也能按台账核。台账编号不替代证据文件，只负责把正文的一句话指回一条可打开的记录。如果编号与文件对不上，以文件为准，并回改正文。

\section{实验与运行总表}

本附录按运行家族给出逐次记录与失败清单。表 B1 覆盖平台四模块的 54 次锁频控制器运行。表 B2 覆盖 M2 除法热点微基准的 6 次运行。表 B3 覆盖收尾的 117 次锁频运行、同机网络的 30 次运行、aarch64 的 12 格与 M1 封存批次，并把失败尝试单列。控制器的“0/完成/是”表示 controller\_exit = 0、reason = worker\_finished、restore\_match = yes。每次进入的 20 个频率探针均在 $\pm 2\%$ 带内且跨度不超过 50000 kHz。

\begin{center}
{\scriptsize
\renewcommand{\arraystretch}{0.95}
\begin{tabularx}{\textwidth}{|c|c|c|c|>{\raggedright\arraybackslash}X|c|c|c|>{\centering\arraybackslash}p{2.1cm}|}
\hline
\apptabhead{序} & \apptabhead{模块} & \apptabhead{运行} & \apptabhead{实现版本} & \apptabhead{类别} & \apptabhead{每类样本} & \apptabhead{日期} & \apptabhead{控制器} & \apptabhead{产物 / 异常} \\
\hline
1 & M5 & s1 31337 CPU2 & 参考 & measure 主类对 & 10000 & 10-16 & 0/完成/是 & CSV+meta / — \\
2 & M5 & s1 31337 CPU2 & 参考 & online dudect & 10000/批 & 10-16 & 0/完成/是 & max t 日志 / — \\
3 & M5 & s2 424242 CPU2 & 参考 & measure 主类对 & 10000 & 10-16 & 0/完成/是 & CSV+meta / — \\
4 & M5 & s2 424242 CPU2 & 参考 & online dudect & 10000/批 & 10-16 & 0/完成/是 & max t 日志 / — \\
5 & M5 & s3 31337 CPU6 & 参考 & measure 主类对 & 10000 & 10-16 & 0/完成/是 & CSV+meta / — \\
6 & M5 & s3 31337 CPU6 & 参考 & online dudect & 10000/批 & 10-16 & 0/完成/是 & max t 日志 / — \\
\hline
7 & M3 & s1 31337 CPU2 & on & measure 主类对 & 10000 & 10-16 & 0/完成/是 & CSV+meta / — \\
8 & M3 & s1 31337 CPU2 & off & measure 主类对 & 10000 & 10-16 & 0/完成/是 & CSV+meta / — \\
9 & M3 & s2 424242 CPU2 & on & measure 主类对 & 10000 & 10-16 & 0/完成/是 & CSV+meta / — \\
10 & M3 & s2 424242 CPU2 & off & measure 主类对 & 10000 & 10-16 & 0/完成/是 & CSV+meta / — \\
11 & M3 & s3 31337 CPU6 & on & measure 主类对 & 10000 & 10-16 & 0/完成/是 & CSV+meta / — \\
12 & M3 & s3 31337 CPU6 & off & measure 主类对 & 10000 & 10-16 & 0/完成/是 & CSV+meta / — \\
13 & M3 & s1 31337 CPU2 & on & online dudect & 10000/批 & 10-16 & 0/完成/是 & max t 日志 / — \\
14 & M3 & s1 31337 CPU2 & off & online dudect & 10000/批 & 10-16 & 0/完成/是 & max t 日志 / — \\
15 & M3 & s2 424242 CPU2 & on & online dudect & 10000/批 & 10-16 & 0/完成/是 & max t 日志 / — \\
16 & M3 & s2 424242 CPU2 & off & online dudect & 10000/批 & 10-16 & 0/完成/是 & max t 日志 / — \\
17 & M3 & s3 31337 CPU6 & on & online dudect & 10000/批 & 10-16 & 0/完成/是 & max t 日志 / — \\
18 & M3 & s3 31337 CPU6 & off & online dudect & 10000/批 & 10-16 & 0/完成/是 & max t 日志 / — \\
\hline

\end{tabularx}
}
\apptabnote{表 B1　平台四模块运行总表（一）。运行编号与正文 4.1、4.3 至 4.7 节一一对应。失败尝试单列于表 B3，不计入正式结果。}
\end{center}

\begin{center}
{\scriptsize
\renewcommand{\arraystretch}{0.95}
\begin{tabularx}{\textwidth}{|c|c|c|c|>{\raggedright\arraybackslash}X|c|c|c|>{\centering\arraybackslash}p{2.1cm}|}
\hline
\apptabhead{序} & \apptabhead{模块} & \apptabhead{运行} & \apptabhead{实现版本} & \apptabhead{类别} & \apptabhead{每类样本} & \apptabhead{日期} & \apptabhead{控制器} & \apptabhead{产物 / 异常} \\
\hline
19 & M4 & s1 31337 CPU2 & 脆弱实现 & measure 主类对 & 10000 & 10-16 & 0/完成/是 & CSV+meta / — \\
20 & M4 & s1 31337 CPU2 & 脆弱实现 & 截断对照 & 10000 & 10-16 & 0/完成/是 & CSV+meta / — \\
21 & M4 & s1 31337 CPU2 & 脆弱实现 & online dudect & 10000/批 & 10-16 & 0/完成/是 & max t 日志 / — \\
22 & M4 & s1 31337 CPU2 & 常数时间实现 & measure 主类对 & 10000 & 10-16 & 0/完成/是 & CSV+meta / — \\
23 & M4 & s1 31337 CPU2 & 常数时间实现 & 截断对照 & 10000 & 10-16 & 0/完成/是 & CSV+meta / — \\
24 & M4 & s1 31337 CPU2 & 常数时间实现 & online dudect & 10000/批 & 10-16 & 0/完成/是 & max t 日志 / — \\
25 & M4 & s2 424242 CPU2 & 脆弱实现 & measure 主类对 & 10000 & 10-16 & 0/完成/是 & CSV+meta / — \\
26 & M4 & s2 424242 CPU2 & 脆弱实现 & 截断对照 & 10000 & 10-16 & 0/完成/是 & CSV+meta / — \\
27 & M4 & s2 424242 CPU2 & 脆弱实现 & online dudect & 10000/批 & 10-16 & 0/完成/是 & max t 日志 / — \\
28 & M4 & s2 424242 CPU2 & 常数时间实现 & measure 主类对 & 10000 & 10-16 & 0/完成/是 & CSV+meta / — \\
29 & M4 & s2 424242 CPU2 & 常数时间实现 & 截断对照 & 10000 & 10-16 & 0/完成/是 & CSV+meta / — \\
30 & M4 & s2 424242 CPU2 & 常数时间实现 & online dudect & 10000/批 & 10-16 & 0/完成/是 & max t 日志 / — \\
31 & M4 & s3 31337 CPU6 & 脆弱实现 & measure 主类对 & 10000 & 10-16 & 0/完成/是 & CSV+meta / — \\
32 & M4 & s3 31337 CPU6 & 脆弱实现 & 截断对照 & 10000 & 10-16 & 0/完成/是 & CSV+meta / — \\
33 & M4 & s3 31337 CPU6 & 脆弱实现 & online dudect & 10000/批 & 10-16 & 0/完成/是 & max t 日志 / — \\
34 & M4 & s3 31337 CPU6 & 常数时间实现 & measure 主类对 & 10000 & 10-16 & 0/完成/是 & CSV+meta / — \\
35 & M4 & s3 31337 CPU6 & 常数时间实现 & 截断对照 & 10000 & 10-16 & 0/完成/是 & CSV+meta / — \\
36 & M4 & s3 31337 CPU6 & 常数时间实现 & online dudect & 10000/批 & 10-16 & 0/完成/是 & max t 日志 / — \\
\hline
37 & M2 & s1 31337 CPU2 & 未修复版 & measure 主类对 & 10000 & 10-16 & 0/完成/是 & CSV+meta / — \\
38 & M2 & s1 31337 CPU2 & 未修复版 & measure 对照类对 & 10000 & 10-16 & 0/完成/是 & CSV+meta / — \\
39 & M2 & s1 31337 CPU2 & 已修复版 & measure 主类对 & 10000 & 10-16 & 0/完成/是 & CSV+meta / — \\
40 & M2 & s1 31337 CPU2 & 已修复版 & measure 对照类对 & 10000 & 10-16 & 0/完成/是 & CSV+meta / — \\
41 & M2 & s1 31337 CPU2 & 未修复版 & online dudect & 10000/批 & 10-16 & 0/完成/是 & max t 日志 / — \\
42 & M2 & s1 31337 CPU2 & 已修复版 & online dudect & 10000/批 & 10-16 & 0/完成/是 & max t 日志 / — \\
43 & M2 & s2 424242 CPU2 & 未修复版 & measure 主类对 & 10000 & 10-16 & 0/完成/是 & CSV+meta / — \\
44 & M2 & s2 424242 CPU2 & 未修复版 & measure 对照类对 & 10000 & 10-16 & 0/完成/是 & CSV+meta / — \\
45 & M2 & s2 424242 CPU2 & 已修复版 & measure 主类对 & 10000 & 10-16 & 0/完成/是 & CSV+meta / — \\
46 & M2 & s2 424242 CPU2 & 已修复版 & measure 对照类对 & 10000 & 10-16 & 0/完成/是 & CSV+meta / — \\
47 & M2 & s2 424242 CPU2 & 未修复版 & online dudect & 10000/批 & 10-16 & 0/完成/是 & max t 日志 / — \\
48 & M2 & s2 424242 CPU2 & 已修复版 & online dudect & 10000/批 & 10-16 & 0/完成/是 & max t 日志 / — \\
49 & M2 & s3 31337 CPU6 & 未修复版 & measure 主类对 & 10000 & 10-16 & 0/完成/是 & CSV+meta / — \\
50 & M2 & s3 31337 CPU6 & 未修复版 & measure 对照类对 & 10000 & 10-16 & 0/完成/是 & CSV+meta / — \\
51 & M2 & s3 31337 CPU6 & 已修复版 & measure 主类对 & 10000 & 10-16 & 0/完成/是 & CSV+meta / — \\
52 & M2 & s3 31337 CPU6 & 已修复版 & measure 对照类对 & 10000 & 10-16 & 0/完成/是 & CSV+meta / — \\
53 & M2 & s3 31337 CPU6 & 未修复版 & online dudect & 10000/批 & 10-16 & 0/完成/是 & max t 日志 / — \\
54 & M2 & s3 31337 CPU6 & 已修复版 & online dudect & 10000/批 & 10-16 & 0/完成/是 & max t 日志 / — \\
\hline
\end{tabularx}
}
\apptabnote{表 B1　平台四模块运行总表（二）。运行编号与正文 4.1、4.3 至 4.7 节一一对应。失败尝试单列于表 B3，不计入正式结果。}
\end{center}

\begin{center}
{\footnotesize
\renewcommand{\arraystretch}{1.0}
\begin{tabularx}{\textwidth}{|c|c|c|c|c|c|>{\raggedright\arraybackslash}X|}
\hline
\apptabhead{序} & \apptabhead{运行} & \apptabhead{实现版本} & \apptabhead{种子 / 核} & \apptabhead{每类样本} & \apptabhead{控制器} & \apptabhead{产物 / 异常} \\
\hline
1 & m2-micro-s1-unfixed & 未修复版 & 31337 / CPU2 & 10000 & 0/完成/是 & CSV+meta / — \\
2 & m2-micro-s1-fixed & 已修复版 & 31337 / CPU2 & 10000 & 0/完成/是 & CSV+meta / — \\
3 & m2-micro-s2-fixed & 已修复版 & 424242 / CPU2 & 10000 & 0/完成/是 & CSV+meta / — \\
4 & m2-micro-s2-unfixed & 未修复版 & 424242 / CPU2 & 10000 & 0/完成/是 & CSV+meta / — \\
5 & m2-micro-s3-unfixed & 未修复版 & 31337 / CPU6 & 10000 & 0/完成/是 & CSV+meta / — \\
6 & m2-micro-s3-fixed & 已修复版 & 31337 / CPU6 & 10000 & 0/完成/是 & CSV+meta / — \\
\hline
\end{tabularx}
}
\apptabnote{表 B2　M2 除法热点微基准（6 次）。窗口只包一次消息编码调用，与完整解封装不可比。6 次运行均通过控制器四判据。}
\end{center}

\begin{center}
{\footnotesize
\renewcommand{\arraystretch}{1.0}
\begin{tabularx}{\textwidth}{|>{\raggedright\arraybackslash}p{2.35cm}|>{\raggedright\arraybackslash}p{4.0cm}|c|c|>{\raggedright\arraybackslash}p{2.55cm}|c|>{\raggedright\arraybackslash}X|}
\hline
\apptabhead{批次} & \apptabhead{运行构成} & \apptabhead{次数} & \apptabhead{日期} & \apptabhead{样本 / 类} & \apptabhead{控制器} & \apptabhead{产物 / 失败尝试} \\
\hline
异常输入对照（格式破损） & 3 种实现 × 3 种子 ×（measure + dudect） & 18 & 10-16 & 10000 & 18/18 & CSV 与 max t 日志。无失败 \\
\hline
性能块 & 9 个实现版本 × 6 个交错块 & 54 & 10-16 & decaps 150000。keygen、encaps 30000。M4 120000 & 54/54 & CSV+meta。18 次无效运行（块号参数错误，退出码 2，无计时值，不计入） \\
\hline
优化级别 & 12 格 × 3 种子 & 36 & 10-16 & 10000 & 36/36 & CSV+meta。1 次被拒（频率探针跨度 54629 kHz，超过 50000 kHz，无计时值） \\
\hline
锁频对照 & 3 种实现 × 3 种子（锁频组） & 9 & 10-16 & 10000 & 9/9 & CSV+meta。不锁频组 9 次另列，不计入 117 \\
\hline
同机网络 & 3 种实现 × 2 次 × 5 档 & 30 & 10-20 & 本地直调 10000。网络各档 2000 & 30/30 & 客户端 RTT 原始值 + 服务端摘要。5 次频率探针拒绝、2 次客户端缺失，均无计时值 \\
\hline
aarch64 格 & 3 种实现 × 4 个优化级别 & 12 & 10-21 & 不采时间 & 无控制器 & Memcheck 日志与反汇编。无失败 \\
\hline
M1 封存批次 & 两个实现版本 × 3 次运行 & 6 & 10-15/16 & 10000 & 6/6 & replication.csv。同一机器与锁频规程，非与 M2 至 M5 同批采集 \\
\hline
\end{tabularx}
}
\apptabnote{表 B3　收尾运行、网络、架构与封存批次。117 = 18 + 54 + 36 + 9，均为锁频控制器运行。不锁频 9 次与失败尝试不并入。}
\end{center}

同机网络场景的逐格结果见表 B4。五档为本地直调、回环、low、mid、high。网络各档每类 2000 个样本，三次运行中的每一次都单独记录。检测预算按 $n = 2(1.96\,\mathrm{SD}/|\Delta|)^2$ 计算，参照效应取本地直调档两次运行的合并值 14.6 ns，只作量级参照。

\begin{center}
{\scriptsize
\renewcommand{\arraystretch}{0.95}
\begin{tabularx}{\textwidth}{|c|c|c|c|c|c|>{\centering\arraybackslash}X|}
\hline
\apptabhead{实现版本} & \apptabhead{次} & \apptabhead{本地直调} & \apptabhead{回环} & \apptabhead{low} & \apptabhead{mid} & \apptabhead{high} \\
\hline
M1 未修复版 & r1 & +37.1 / 可分辨 & +63.1 / 未检出 & $-801.3$ / 未检出 & +6515.5 / 未检出 & +39569.7 / 未检出 \\
M1 未修复版 & r2 & $-7.9$ / 未检出 & $-12.4$ / 未检出 & +263.0 / 未检出 & $-1863.5$ / 未检出 & $-51509.7$ / 未检出 \\
M1 已修复版 & r1 & +11.9 / 未检出 & $-52.3$ / 未检出 & +736.4 / 未检出 & $-3151.9$ / 未检出 & +72863.2 / 未检出 \\
M1 已修复版 & r2 & +7.0 / 未检出 & +5.9 / 未检出 & $-942.0$ / 未检出 & $-6885.7$ / 未检出 & +39109.8 / 未检出 \\
M5 参考 & r1 & $-11.3$ / 未检出 & $-52.8$ / 未检出 & $-665.1$ / 未检出 & +2861.4 / 未检出 & $-30649.4$ / 未检出 \\
M5 参考 & r2 & +2.7 / 未检出 & $-57.7$ / 未检出 & $-455.4$ / 未检出 & +654.6 / 未检出 & $-61629.6$ / 未检出 \\
\hline
\multicolumn{2}{|l|}{分辨 14.6 ns 所需样本 / 类} & 13098 & 38417 & 64636555 & 1003153317 & 97586175893 \\
\multicolumn{2}{|l|}{本轮实际样本 / 类} & 10000 & 2000 & 2000 & 2000 & 2000 \\
\hline
\end{tabularx}
}
\apptabnote{表 B4　同机网络五档的均值差（ns）与判定。网络各档在本预算下未检出。这是同机模拟、非真实远程，所需样本比实际高 1.3 至 7.7 个数量级。}
\end{center}

\begin{center}
{\footnotesize
\renewcommand{\arraystretch}{1.0}
\begin{tabularx}{\textwidth}{|>{\raggedright\arraybackslash}p{2.1cm}|c|c|c|c|>{\raggedright\arraybackslash}p{3.6cm}|>{\raggedright\arraybackslash}X|}
\hline
\apptabhead{实现版本} & \apptabhead{-O0} & \apptabhead{-O2} & \apptabhead{-O3} & \apptabhead{-Os} & \apptabhead{反汇编形态} & \apptabhead{判定} \\
\hline
M1 未修复版 & 0 / 0 & 0 / 0 & 0 / 0 & 0 / 0 & \texttt{sbfx}、\texttt{and}、\texttt{csel} & 四格均无秘密相关分支 \\
M1 已修复版 & 0 / 0 & 0 / 0 & 0 / 0 & 0 / 0 & 掩码选择 & 四格均无命中 \\
M5 参考 & 0 / 0 & 0 / 0 & 0 / 0 & 0 / 0 & 掩码选择 + 值屏障 & 四格均无命中 \\
x86\_64 M1 未修复版（对照） & 0 / 0 & 512 / 8 & 512 / 15 & 512 / 1 & \texttt{bt}/\texttt{jae} 或逐位测试分支 & x86\_64 三格有命中 \\
\hline
\end{tabularx}
}
\apptabnote{表 B5　aarch64 12 格与 x86\_64 同格对照（错误数 / 上下文数）。}
\end{center}

表 B1 至表 B5 的共同口径是：一条运行记录只对应一个模块、一个实现版本、一个种子与一个控制器终态。同一模块的 measure 与 dudect 分开记，不合并成一次。54 次平台家族、6 次微基准与 117 次收尾运行合计 177 次控制器运行。M1 的 6 次主类对属于封存批次，不并入 177。失败尝试单列，不删、不改名、不重算。它们没有计时值，因此既不构成阳性也不构成阴性。表 B4 的网络与表 B5 的 aarch64 不进任何 Holm 家族，时间证据不与主平台合并。

样本预算按主类对设计：每类 10000 次、三次独立运行，换种子与换物理核。微基准只包一个热点函数，与完整解封装不可比。网络档每类 2000 次，检测预算按同一式子反推所需样本。预算不同的表不能横向比较 Δ 的大小，只能各自回答“在本预算下是否可分辨”。这也是表 B3 把 54、6、117 分开列、不合成一个总数的原因。

日期与批次也要一起读。平台四模块、微基准、性能、优化与锁频对照集中在 10 月 16 日。异常输入对照同日。同机网络在 10 月 20 日。aarch64 在 10 月 21 日。M1 的主类对是 10 月 15 日至 16 日的封存批次。同一台机器、同一锁频规程，不等于同一批采集。跨日数据只做同口径对照，不把日期差当成处理效应。表 B1 的 54 次运行全部有控制器终态与原始 CSV，表 B3 的 117 次运行另有逐份 terminal 记录。两类记录都保留失败尝试，避免“只留成功”造成的工作量虚高。

失败尝试单列还有一层作用：它们能区分“设备没跑起来”与“跑了但没有差异”。控制器拒绝、参数错误、客户端缺失都属于前者，没有计时值。它们既不支持阳性，也不支持阴性。把它们写成“失败”而不是删掉，读者才能判断 54、6、117 这些数字是不是从全部尝试里挑出来的。本附录的所有计数都以“成功且四判据通过”为分子，以“成功 + 失败尝试”为分母。分母写在表里，分子写在表注里。

\section{统计公式与复算说明}

本附录给出统计层用到的七个公式、本次参数与手算例。公式与统计电池、固定效应合并工具的实现一致。dudect 与 RTLF 的工具行为按上游实现\cite{dudect,rtlf}。手算例只用合成数核对公式，不替代正式数据。正式结果见附录 D 与正文 4.3 至 4.7 节。

\paragraph{（1）Welch 检验与自由度。}
\begin{equation}\tag{C1}
t=\frac{\Delta}{\sqrt{s_0^2/n_0+s_1^2/n_1}},\qquad
\mathrm{df}=\frac{\left(s_0^2/n_0+s_1^2/n_1\right)^2}{\dfrac{(s_0^2/n_0)^2}{n_0-1}+\dfrac{(s_1^2/n_1)^2}{n_1-1}}.
\end{equation}
其中 $\Delta=\bar x_0-\bar x_1$，$s_0^2$、$s_1^2$ 为样本方差，$n_0=n_1=10000$。本次不称其为有限样本精确 $t$ 区间。$p$ 用 $\mathrm{erfc}(|t|/\sqrt2)$ 的正态近似。

\paragraph{（2）95\% 置信区间。}
\begin{equation}\tag{C2}
\left[\Delta-z_{0.975}\,\mathrm{SE},\ \Delta+z_{0.975}\,\mathrm{SE}\right],\qquad
\mathrm{SE}=\sqrt{s_0^2/n_0+s_1^2/n_1},\quad z_{0.975}=1.9599639845.
\end{equation}
区间含 0 记“未检出”。区间不含 0 记“可区分”。未检出必须同时给出 MDE。

\paragraph{（3）Holm 逐步校正。}
\begin{equation}\tag{C3}
p_{(1)}\le p_{(2)}\le\cdots\le p_{(m)},\qquad
\text{第 }i\text{ 项阈值}=\frac{\alpha}{m-i+1}.
\end{equation}
从最小 $p$ 开始逐项比较。第一个不满足即停止，其后各项均不拒绝。全族 $\alpha=0.05$。M1、M2、M3、M4 的 $m=6$，M5 的 $m=3$。

\paragraph{（4）逆方差固定效应合并。}
\begin{equation}\tag{C4}
w_i=\frac{1}{\mathrm{SE}_i^2},\qquad
\Delta_{\mathrm{pool}}=\frac{\sum_i w_i\Delta_i}{\sum_i w_i},\qquad
\mathrm{SE}_{\mathrm{pool}}=\frac{1}{\sqrt{\sum_i w_i}}.
\end{equation}
每次运行的 $\mathrm{SE}_i$ 由该次 95\% 区间反推：$\mathrm{SE}_i=(\text{上界}-\text{下界})/(2z_{0.975})$。

\paragraph{（5）异质性 $I^2$。}
\begin{equation}\tag{C5}
Q=\sum_i w_i\left(\Delta_i-\Delta_{\mathrm{pool}}\right)^2,\qquad
I^2=\max\left(0,\frac{Q-(k-1)}{Q}\right).
\end{equation}
$k$ 为成员数。$I^2$ 高时，合并区间只作描述，判定仍以预注册的 Holm 结果为准。

\paragraph{（6）最小可检测效应。}
\begin{equation}\tag{C6}
\mathrm{MDE}=\left(z_{1-\alpha/2}+z_{1-\beta}\right)\sqrt{\frac{s_0^2}{n_0}+\frac{s_1^2}{n_1}},
\end{equation}
取 $\alpha=0.05$、$1-\beta=0.80$，$z_{0.975}=1.959963985$、$z_{0.80}=0.841621234$。Holm 口径把 $\alpha$ 换成 $\alpha_{\mathrm{eff}}=\alpha/m$。M1、M2、M3、M4 的 $z_{1-\alpha_{\mathrm{eff}}/2}=2.638257273$，M5 为 $2.393979800$。MDE 是事后功效描述，不改变已提交的判定。

\paragraph{（7）RTLF 秩分位数 bootstrap。}
RTLF 对两类样本的分位数差做 bootstrap 重采样，$\alpha=0.01$，固定种子 2026101605，重采样 10000 次，输出二值结论：退出码 10 为未检出、11 为检出。它不产生 $p$ 值，不进入 Holm 家族，也不与均值检验互相替代。

\begin{center}
{\footnotesize
\renewcommand{\arraystretch}{1.05}
\begin{tabularx}{\textwidth}{|c|>{\raggedright\arraybackslash}p{3.2cm}|>{\raggedright\arraybackslash}p{5.6cm}|>{\raggedright\arraybackslash}X|}
\hline
\apptabhead{公式} & \apptabhead{符号} & \apptabhead{本次取值} & \apptabhead{用途} \\
\hline
C1 & $t$、$\mathrm{df}$、$s^2$、$n$ & $n_0=n_1=10000$。$\alpha=0.05$ & 单次运行的均值差检验与自由度 \\
\hline
C2 & $\mathrm{SE}$、$z_{0.975}$ & $z_{0.975}=1.9599639845$ & 每次运行的 95\% 区间 \\
\hline
C3 & $p_{(i)}$、$m$、$\alpha$ & $m=6$（M5 为 3）。$\alpha=0.05$ & 家族内多重比较 \\
\hline
C4 & $w_i$、$\Delta_{\mathrm{pool}}$ & 三次运行按区间反推权重 & 跨运行合并 \\
\hline
C5 & $Q$、$k$ & $k=6$（M5 为 3） & 异质性描述 \\
\hline
C6 & MDE、$z_{1-\beta}$ & $z_{0.80}=0.841621234$。Holm 口径 $\alpha_{\mathrm{eff}}=\alpha/m$ & 量化未检出 \\
\hline
RTLF & 分位数、bootstrap & $\alpha=0.01$。种子 2026101605。10000 次 & 族外分布交叉 \\
\hline
\end{tabularx}
}
\apptabnote{表 C1　统计公式与参数。公式与统计电池、固定效应合并工具的实现一致。手算例只用于核对，不替代正式数据。}
\end{center}

\paragraph{手算例一（Welch 与区间）。}
取两组各 3 个数：$x_0=\{10,12,11\}$，$x_1=\{7,9,8\}$。均值 $\bar x_0=11$、$\bar x_1=8$，$\Delta=3$。样本方差 $s_0^2=s_1^2=1$。$\mathrm{SE}=\sqrt{1/3+1/3}=0.8165$，$t=3.674$，$\mathrm{df}=4$，双侧 $p\approx0.021$。95\% 区间为 $3\pm1.96\times0.8165=[1.40,4.60]$，不含 0。这个例子只说明式 C1、C2 的算法，不构成任何实验结论。

\paragraph{手算例二（Holm）。}
取 $m=4$ 与 $p=\{0.001,0.020,0.030,0.040\}$。阈值依次为 $0.0125$、$0.0167$、$0.025$、$0.05$。第一项 $0.001<0.0125$，拒绝。第二项 $0.020>0.0167$，停止。其后两项均不拒绝。Holm 的停止规则使“前面不拒绝”自动覆盖后面，不需要逐项再比。

\paragraph{手算例三（合并与 $I^2$）。}
三次运行的 $\Delta=\{10,14,8\}$，$\mathrm{SE}=\{2,2,4\}$。

权重 $w=\{0.25,0.25,0.0625\}$，$\sum w=0.5625$。\par
$\Delta_{\mathrm{pool}}=(2.5+3.5+0.5)/0.5625=11.556$。$\mathrm{SE}_{\mathrm{pool}}=1/\sqrt{0.5625}=1.333$。95\% 区间 $11.556\pm2.613=[8.94,14.17]$。$Q=0.25(1.556)^2+0.25(2.444)^2+0.0625(3.556)^2=2.89$，$k-1=2$，$I^2=(2.89-2)/2.89=30.8\%$。

\paragraph{手算例四（MDE）。}
取 $s_0=700$ ns、$s_1=600$ ns、$n_0=n_1=10000$、$\alpha=0.05$、功效 0.80。$\mathrm{MDE}=2.801585\times\sqrt{700^2/10000+600^2/10000}=2.801585\times9.2195=25.83$ ns。若改用 $m=6$ 的 Holm 口径，$z_{1-\alpha_{\mathrm{eff}}/2}=2.638257$，MDE 变为 $(2.638257+0.841621)\times9.2195=32.08$ ns。

\paragraph{口径说明。}
M1 族在采集后发生过一次成员调整：dudect 的多批最大 $t$ 不是有效 $p$ 值，因此从参数族移出，原始记录另存。两种口径下的家族级判定一致，正文与附录 D 的判定均未改变。这里不展开成员数，避免把一次口径调整写成新的实验事实。

区间与 $p$ 值采用大样本正态近似；MDE 是基于实测标准差的事后功效描述，不改变预注册判定。异质性较高时，合并区间只作描述；M4 的家族合并混合了脆弱实现与常数时间实现，需结合分版本结果解释。RTLF 只给二值结论，不转成 $p$ 值，也不与均值检验合并。

Holm 与 MDE 是两件事：Holm 控制家族错误率，MDE 描述可检测性。一个阴性结论要同时给两者，读者才能判断“没拒绝”是因为效应小，还是因为样本不够。手算例只用来核对公式，不替代正式数据。正式数据在附录 D，逐次值与家族判定并列。

\section{模块级完整结果}

本附录给出 M1 至 M5 的逐次运行结果。$\Delta$ 为 class0 减 class1 的均值差。95\% 区间为正态近似。$p$ 为双侧原始值。Holm 列为该成员在预注册家族中的判定。dudect 与 RTLF 为族外交叉，不进入 Holm。M1 的逐次 $p$ 由该次 $\Delta$ 与区间反推 $\mathrm{SE}$ 后按式 C1、C2 复算，属 B 类。其余 $p$ 取自成员表原件。M4 的 $p=0$ 是双精度下溢，不写“精确为零”。

\begin{center}
{\scriptsize
\renewcommand{\arraystretch}{0.95}
\begin{tabularx}{\textwidth}{|c|c|c|c|c|c|c|c|>{\raggedright\arraybackslash}X|}
\hline
\apptabhead{实现版本} & \apptabhead{运行} & \apptabhead{$\Delta$ / ns} & \apptabhead{95\% CI / ns} & \apptabhead{$p$} & \apptabhead{Holm} & \apptabhead{dudect $\max|t|$} & \apptabhead{RTLF} & \apptabhead{MDE / ns} \\
\hline
未修复版 & s1 31337 CPU2 & +39.337 & [21.658, 57.017] & $1.30\times10^{-5}$ & 族内拒绝 & 26.41 & 检出 & 25.27 / 31.39 \\
未修复版 & s2 424242 CPU2 & +31.117 & [6.549, 55.684] & 0.0131 & 族内拒绝 & 40.64 & 检出 & 25.27 / 31.39 \\
未修复版 & s3 31337 CPU6 & +22.337 & [0.282, 44.392] & 0.0471 & 族内拒绝 & 14.92 & 检出 & 25.27 / 31.39 \\
已修复版 & s1 31337 CPU2 & $-13.209$ & [$-33.238$, 6.820] & 0.1962 & 不拒绝 & 1.15 & 未检出 & 28.63 / 35.56 \\
已修复版 & s2 424242 CPU2 & +8.354 & [$-52.988$, 69.696] & 0.7895 & 不拒绝 & 1.12 & 未检出 & 28.63 / 35.56 \\
已修复版 & s3 31337 CPU6 & $-4.885$ & [$-24.811$, 15.042] & 0.6309 & 不拒绝 & 1.31 & 未检出 & 28.63 / 35.56 \\
\hline
\multicolumn{2}{|l|}{合并（未修复版）} & +32.311 & [20.282, 44.339] & $1.30\times10^{-5}$ & 拒绝 & — & — & 25.27 / 31.39 \\
\multicolumn{2}{|l|}{合并（已修复版）} & $-8.150$ & [$-21.916$, 5.616] & 0.196 & 不拒绝 & — & — & 28.63 / 35.56 \\
\hline
\end{tabularx}
}
\apptabnote{表 D1　M1 逐次结果。逐次值取自封存原始数据。“未检出”附 MDE。M1 的家族主检验 $p=1.30\times10^{-5}$、阈值 0.00833。已修复版主检验 $p=0.196$、阈值 0.05（末位成员）后不拒绝。}
\end{center}

\begin{center}
{\scriptsize
\renewcommand{\arraystretch}{0.95}
\begin{tabularx}{\textwidth}{|c|c|c|c|c|c|c|c|>{\raggedright\arraybackslash}X|}
\hline
\apptabhead{实现版本} & \apptabhead{运行} & \apptabhead{$\Delta$ / ns} & \apptabhead{95\% CI / ns} & \apptabhead{$p$} & \apptabhead{Holm} & \apptabhead{dudect $\max|t|$} & \apptabhead{RTLF} & \apptabhead{MDE / ns} \\
\hline
未修复版 & r1 31337 CPU2 & +10.530 & [$-6.396$, 27.455] & 0.223 & 不拒绝 & 4.68 & 检出 & 19.9 至 54.0 / 24.7 至 67.1 \\
未修复版 & r2 424242 CPU2 & +19.126 & [0.028, 38.224] & 0.0497 & 不拒绝 & 3.05 & 检出 & 19.9 至 54.0 / 24.7 至 67.1 \\
未修复版 & r3 31337 CPU6 & +9.983 & [$-3.915$, 23.881] & 0.159 & 不拒绝 & 3.66 & 检出 & 19.9 至 54.0 / 24.7 至 67.1 \\
已修复版 & r1 31337 CPU2 & +21.657 & [$-16.123$, 59.436] & 0.261 & 不拒绝 & 36.27 & 检出 & 19.9 至 54.0 / 24.7 至 67.1 \\
已修复版 & r2 424242 CPU2 & $-17.080$ & [$-45.991$, 11.830] & 0.247 & 不拒绝 & 30.94 & 检出 & 19.9 至 54.0 / 24.7 至 67.1 \\
已修复版 & r3 31337 CPU6 & $-15.512$ & [$-31.925$, 0.900] & 0.0640 & 不拒绝 & 61.91 & 检出 & 19.9 至 54.0 / 24.7 至 67.1 \\
\hline
\multicolumn{2}{|l|}{合并} & +4.58 & [$-3.08$, 12.25] & — & 拒绝数 0 & — & — & 24.7 至 67.1 \\
\hline
\end{tabularx}
}
\apptabnote{表 D2　M2 逐次结果。逐次值由封存原始数据重列。“未检出”附 MDE。dudect 阳性落在已修复版、RTLF 两个版本均检出，均不可归因于除法变体，只作族外登记。}
\end{center}

\begin{center}
{\scriptsize
\renewcommand{\arraystretch}{0.95}
\begin{tabularx}{\textwidth}{|c|c|c|c|c|c|c|c|>{\raggedright\arraybackslash}X|}
\hline
\apptabhead{实现版本} & \apptabhead{运行} & \apptabhead{$\Delta$ / ns} & \apptabhead{95\% CI / ns} & \apptabhead{$p$} & \apptabhead{Holm} & \apptabhead{dudect $\max|t|$} & \apptabhead{RTLF} & \apptabhead{MDE / ns} \\
\hline
on & r1 31337 CPU2 & $-10.723$ & [$-29.703$, 8.258] & 0.268 & 不拒绝 & 1.93 & 未检出 & 14.2 至 27.1 / 17.7 至 33.7 \\
on & r2 424242 CPU2 & $-5.763$ & [$-17.016$, 5.489] & 0.315 & 不拒绝 & 3.41 & 未检出 & 14.2 至 27.1 / 17.7 至 33.7 \\
on & r3 31337 CPU6 & +7.126 & [$-3.497$, 17.749] & 0.189 & 不拒绝 & 2.81 & 未检出 & 14.2 至 27.1 / 17.7 至 33.7 \\
off & r1 31337 CPU2 & $-5.464$ & [$-20.612$, 9.683] & 0.480 & 不拒绝 & 2.24 & 未检出 & 14.2 至 27.1 / 17.7 至 33.7 \\
off & r2 424242 CPU2 & +4.800 & [$-5.816$, 15.417] & 0.375 & 不拒绝 & 3.37 & 未检出 & 14.2 至 27.1 / 17.7 至 33.7 \\
off & r3 31337 CPU6 & $-3.775$ & [$-13.732$, 6.183] & 0.457 & 不拒绝 & 3.03 & 未检出 & 14.2 至 27.1 / 17.7 至 33.7 \\
\hline
\multicolumn{2}{|l|}{合并} & $-0.73$ & [$-5.56$, 4.10] & — & 拒绝数 0 & — & — & 17.7 至 33.7 \\
\hline
\end{tabularx}
}
\apptabnote{表 D3　M3 逐次结果。逐次值来自封存记录。“未检出”附 MDE。M3 为人工构造校准，不是上游缺陷。}
\end{center}

\begin{center}
{\scriptsize
\renewcommand{\arraystretch}{0.95}
\setlength{\tabcolsep}{2.5pt}
\begin{tabularx}{\textwidth}{|c|c|c|c|c|c|c|c|>{\raggedright\arraybackslash}X|}
\hline
\apptabhead{实现版本} & \apptabhead{运行} & \apptabhead{$\Delta$ / ns} & \apptabhead{95\% CI / ns} & \apptabhead{$p$} & \apptabhead{Holm} & \apptabhead{dudect $\max|t|$} & \apptabhead{RTLF} & \apptabhead{MDE / ns} \\
\hline
脆弱实现 & r1 31337 CPU2 & $-4.102$ & [$-4.232$, $-3.973$] & 下溢 0 & 拒绝 & 356.6 & 检出 & 0.19 至 1.19 / 1.47 \\
脆弱实现 & r2 424242 CPU2 & $-4.351$ & [$-5.181$, $-3.521$] & \mbox{$9.23\times10^{-25}$} & 拒绝 & 396.9 & 检出 & 0.19 至 1.19 / 1.47 \\
脆弱实现 & r3 31337 CPU6 & $-3.515$ & [$-4.226$, $-2.805$] & \mbox{$3.14\times10^{-22}$} & 拒绝 & 395.6 & 检出 & 0.19 至 1.19 / 1.47 \\
常数时间实现 & r1 31337 CPU2 & +0.415 & [$-0.523$, 1.353] & 0.386 & 不拒绝 & 2.06 & 未检出 & 0.19 至 1.19 / 1.47 \\
常数时间实现 & r2 424242 CPU2 & $-0.227$ & [$-0.800$, 0.346] & 0.438 & 不拒绝 & 1.90 & 未检出 & 0.19 至 1.19 / 1.47 \\
常数时间实现 & r3 31337 CPU6 & +0.347 & [$-0.460$, 1.155] & 0.399 & 不拒绝 & 2.27 & 未检出 & 0.19 至 1.19 / 1.47 \\
\hline
\multicolumn{2}{|l|}{合并（脆弱实现）} & $-4.089$ & [$-4.215$, $-3.963$] & — & 拒绝 3 项 & — & — & 1.47（最坏） \\
\multicolumn{2}{|l|}{合并（常数时间实现）} & +0.055 & [$-0.363$, 0.473] & — & 不拒绝 & — & — & 1.47（最坏） \\
\hline
\end{tabularx}
}
\apptabnote{表 D4　M4 分版本逐次结果。逐次值由原始成员表复列。“未检出”附 MDE。$p=0$ 是双精度下溢，不写“精确为零”。}
\end{center}

\begin{center}
{\scriptsize
\renewcommand{\arraystretch}{0.95}
\begin{tabularx}{\textwidth}{|c|c|c|c|c|c|c|c|>{\raggedright\arraybackslash}X|}
\hline
\apptabhead{实现版本} & \apptabhead{运行} & \apptabhead{$\Delta$ / ns} & \apptabhead{95\% CI / ns} & \apptabhead{$p$} & \apptabhead{Holm} & \apptabhead{dudect $\max|t|$} & \apptabhead{RTLF} & \apptabhead{MDE / ns} \\
\hline
参考 & r1 31337 CPU2 & $-5.596$ & [$-18.912$, 7.720] & 0.410 & 不拒绝 & 2.10 & 未检出 & 11.7 至 24.5 / 13.5 至 28.3 \\
参考 & r2 424242 CPU2 & $-6.946$ & [$-24.103$, 10.211] & 0.427 & 不拒绝 & 1.98 & 未检出 & 11.7 至 24.5 / 13.5 至 28.3 \\
参考 & r3 31337 CPU6 & $-4.819$ & [$-13.019$, 3.382] & 0.249 & 不拒绝 & 2.82 & 未检出 & 11.7 至 24.5 / 13.5 至 28.3 \\
\hline
\multicolumn{2}{|l|}{合并} & $-5.30$ & [$-11.77$, 1.16] & — & 拒绝数 0 & — & — & 13.5 至 28.3 \\
\hline
\end{tabularx}
}
\apptabnote{表 D5　M5 逐次结果。逐次值取自封存记录；“未检出”附 MDE。}
\end{center}

逐次表的读法与正文不同：正文给分层结论，附录给每次运行的原始量。Δ 的符号由 class0 减 class1 决定，同一模块内方向不一致时，不取绝对值平均。先看每次区间是否含 0，再看 Holm 家族是否拒绝。M2 的未修复版描述性合并为正、已修复版为负，这是实现版本间的漂移，不是泄漏降低。M4 的家族合并被两个版本混合拉宽，按实现版本分开后脆弱实现与常数时间实现的区间才可读。M1 的逐次 p 是 B 类复算值，版本合并 p 才是族内判定。两者并列，不互相替代。

五个模块呈现不同的证据组合：M1 定位与计时均阳性，官方修复后计时未检出；M2 定位阳性而端到端与微基准未检出；M3 为人工校准；M4 为定位和计时均阳性的正对照，常数时间实现未检出；M5 为阴性对照，七种工具均未报阳性。逐次值与模块结论分别见表 D1--D5。

MDE 也要逐模块读。M1 未修复版的检测预算约 25 至 31 ns，已修复版约 29 至 36 ns。M2 为 25 至 67 ns。M3 为 18 至 34 ns。M4 最坏 1.47 ns。M5 为 14 至 28 ns。观测效应小于本模块 MDE 时写“未检出”，不写“零效应”。M4 的 MDE 小，是因为比较窗口只有几十纳秒、方差也小。这不表示 M4 的结论可以外推到完整解封装。反过来，M2 的 MDE 大，除法热点即使有几十纳秒的效应也可能落在预算之外，所以定位阳性与计时阴性可以同时成立。

逐次 p 的来源分两类：M2 至 M5 直接取自成员表原件。M1 的逐次 p 由该次 Δ 与区间反推 SE 后按正态近似复算。复算值只用于附录对照，族内判定仍用版本合并 p。把两类值并列，是为了让读者看到“逐次方向一致”与“家族拒绝”不是同一个层次。前者描述三次运行，后者描述预注册检验。

\section{定位工具输出摘要}

本附录只摘录各工具的关键输出行，不贴完整日志。表 E1 覆盖统计与动态标记工具，表 E2 覆盖形式化、访存与反汇编工具。“未运行”按原因登记，不填阴性。工具版本写在各行括号内。同一工具在不同模块上的结论不可互相替代。

\begin{center}
{\scriptsize
\renewcommand{\arraystretch}{0.9}
\begin{tabularx}{\textwidth}{|>{\raggedright\arraybackslash}p{2.55cm}|c|>{\raggedright\arraybackslash}X|>{\raggedright\arraybackslash}p{2.0cm}|>{\raggedright\arraybackslash}p{2.55cm}|>{\raggedright\arraybackslash}p{2.35cm}|}
\hline
\apptabhead{工具（版本）} & \apptabhead{对象} & \apptabhead{关键输出} & \apptabhead{上下文 / 样本} & \apptabhead{位置 / 地址} & \apptabhead{结论} \\
\hline
dudect（上游头文件快照） & M1 & 未修复版 $\max|t|$ 26.41 / 40.64 / 14.92。已修复版 1.15 / 1.12 / 1.31 & 3 次运行，21 批 & 完整解封装 & 检出（未修复版） \\
dudect & M2 & 未修复版 4.68 / 3.05 / 3.66。已修复版 36.27 / 30.94 / 61.91 & 3 次运行 & 完整解封装 & 未检出。阳性落在已修复版，不可归因 \\
dudect & M3 & on 1.93 / 3.41 / 2.81。off 2.24 / 3.37 / 3.03 & 3 次运行 & 完整解封装 & 未检出 \\
dudect & M4 & 脆弱实现 356.6 / 396.9 / 395.6。常数时间实现 2.06 / 1.90 / 2.27 & 3 次运行 & 32 B 比较调用 & 检出（脆弱实现） \\
dudect & M5 & 2.10 / 1.98 / 2.82 & 3 次运行 & 完整解封装 & 未检出 \\
\hline
RTLF（R 4.3.3，$\alpha=0.01$） & M1 & 未修复版 2/2 exit 11。已修复版 2/2 exit 10 & 2 次运行 & 完整解封装 & 检出（未修复版） \\
RTLF & M2 & 主类对 6/6 exit 11（两个版本各 3）。对照 6/6 exit 10 & 6 个数据集 & 完整解封装 & 检出。两个版本均检出，不可归因 \\
RTLF & M3 & 主类对与对照共 12 个数据集全部 exit 10 & 12 个数据集 & 完整解封装 & 未检出 \\
RTLF & M4 & 脆弱实现 6/6 exit 11（含截断）。常数时间实现 6/6 exit 10 & 12 个数据集 & 32 B 比较调用 & 检出（脆弱实现） \\
RTLF & M5 & 主类对与对照共 6/6 exit 10 & 6 个数据集 & 完整解封装 & 未检出 \\
\hline
Welch + Holm（统计电池） & M1 & 主检验 $p=1.30\times10^{-5}$，阈值 0.00833。已修复版 $p=0.196$ & 每类 10000 & 完整解封装 & 检出（未修复版） \\
Welch + Holm & M2 & 合并 $\Delta=+4.58$ ns，CI $[-3.08,12.25]$。拒绝 0 & 6 项 & 完整解封装 & 未检出 \\
Welch + Holm & M3 & 合并 $\Delta=-0.73$ ns，CI $[-5.56,4.10]$。拒绝 0 & 6 项 & 完整解封装 & 未检出 \\
Welch + Holm & M4 & 脆弱实现 3 项拒绝。常数时间实现 0 项 & 6 项 & 32 B 比较调用 & 检出（脆弱实现） \\
Welch + Holm & M5 & 合并 $\Delta=-5.30$ ns，CI $[-11.77,1.16]$。拒绝 0 & 3 项 & 完整解封装 & 未检出 \\
\hline
Memcheck（Valgrind 3.22.0） & M1 & 未修复版 1 error / 1 context。已修复版 0 / 0 & 完整解封装 & \texttt{消息编码函数第 179 行} & 检出（未修复版） \\
Memcheck & M2 & 未运行（完整解封装未跑。仅关 varlat 对照 0/0） & — & — & 未运行，不填阴性 \\
Memcheck & M3 & on 1280 / 2。off 0 / 0 & 10 次完整解封装 & \texttt{消息编码函数第 714 行} & 检出（人工校准） \\
Memcheck & M4 & 32 / 1。常数时间实现 0 / 0 & 完整调用 & \texttt{比较内核第 9 行} & 检出（人工校准） \\
Memcheck & M5 & 三模式各 10 次，0 / 0 & 30 次完整解封装 & 完整调用链 & 未检出 \\
\hline
varlat（Valgrind 3.25.1） & M1 & 未运行（机制不在命中类别） & — & — & 未运行 \\
varlat & M2 & 未修复版 1280 / 3。已修复版 0 / 0 & 3 个热点 & 消息解码、压缩、向量压缩 & 检出（未修复版） \\
varlat & M3 & 未运行 & — & — & 未运行 \\
varlat & M4 & 未运行 & — & — & 未运行 \\
varlat & M5 & varlat=yes 0 / 0 & 10 次完整解封装 & 完整调用链 & 未检出 \\
\hline
\end{tabularx}
}
\apptabnote{表 E1　统计与动态标记工具输出摘要。“未运行”表示该单元没有执行记录。}
\end{center}

\begin{center}
{\scriptsize
\renewcommand{\arraystretch}{0.95}
\begin{tabularx}{\textwidth}{|>{\raggedright\arraybackslash}p{2.55cm}|c|>{\raggedright\arraybackslash}X|>{\raggedright\arraybackslash}p{2.0cm}|>{\raggedright\arraybackslash}p{2.55cm}|>{\raggedright\arraybackslash}p{2.35cm}|}
\hline
\apptabhead{工具（版本）} & \apptabhead{对象} & \apptabhead{关键输出} & \apptabhead{上下文 / 样本} & \apptabhead{位置 / 地址} & \apptabhead{结论} \\
\hline
Binsec（0.11.3 + Z3 4.15.3） & M1 & 未运行（无对应抽取内核） & — & — & 未运行 \\
Binsec & M2 & 未运行（$div$ 类来源不在已开检查项） & — & — & 未运行 \\
Binsec & M3 & 未运行 & — & — & 未运行 \\
Binsec & M4 & 脆弱实现 \texttt{insecure}，反例 \texttt{0x8049047}。常数时间实现 \texttt{bounded secure} & 32/32 与 16/16 路径 & \texttt{比较内核第 9 行} & 检出（人工校准） \\
Binsec & M5 & 重加密比较 1088/1088、消息编码 1/1 路径完成，无秘密相关控制流 & 路径完整 & 重加密比较与消息编码 & 未检出（有界） \\
\hline
Microwalk（源码自建 + Pin 3.28） & M1 & 未运行（官方镜像不可用，Pin 缺失） & — & — & 未运行 \\
Microwalk & M2 & 未运行 & — & — & 未运行 \\
Microwalk & M3 & 互信息 1.922 bit（全场最高）。off 0.722 bit & 5 个公开用例 & 两处秘密下标 load & 检出（人工校准） \\
Microwalk & M4 & 未运行 & — & — & 未运行 \\
Microwalk & M5 & 未运行 & — & — & 未运行 \\
\hline
Mona（上游脚本） & M1 & 未修复版无分离。已修复版 2/6 出现箱体分离 & 6 份输入 & 完整解封装 & 未检出。方向不符 \\
Mona & M2 & 未修复版无分离。s3 已修复版 1/8 出现分离 & 6 批 & 完整解封装 & 未检出。方向不符 \\
Mona & M3 & 未运行（未接入） & — & — & 未运行 \\
Mona & M4 & 未运行（未接入） & — & — & 未运行 \\
Mona & M5 & 未运行（未接入） & — & — & 未运行 \\
\hline
反汇编（发行版 objdump） & M1 & \texttt{bt}+\texttt{jae} → \texttt{bt}+\texttt{setb}+\texttt{call cmov\_int16} & 两版产物 & 消息编码函数 & 检出（未修复版） \\
反汇编 & M2 & 三热点各 1 条 \texttt{div} → 0。消息编码 \texttt{bt}/\texttt{jae} 两个版本中的指令逐字节相同 & 两版产物 & 消息解码、压缩、向量压缩 & 检出。残留分支 \\
反汇编 & M3 & 两条 \texttt{movzwl} 秘密下标 load & on/off 两版 & 静态偏移 \texttt{0x5CED}/\texttt{0x5D15} & 检出（人工校准） \\
反汇编 & M4 & 3 条条件跳转中 1 条数据相关 & 脆弱实现/常数时间实现 & 比较内核 & 检出（人工校准） \\
反汇编 & M5 & 无秘密相关分支。消息编码唯一跳转是循环回边。sk 检查为公开数据 & 完整产物 & 完整调用链 & 未发现 \\
\hline
\end{tabularx}
}
\apptabnote{表 E2　形式化、访存与反汇编工具输出摘要。“未运行”表示该单元没有执行记录。}
\end{center}

45 格中 31 格有真实产物，14 格未运行。未运行的原因包括工具未接入该模块、缺少运行依赖或该模块没有对应的抽取内核。不同工具的判据可能给出不同读法：RTLF 在 M2 两个版本均报告分布差异，dudect 的阳性落在已修复版，预注册 Welch+Holm 对两个版本均未拒绝。综合时还应结合 M2 与 M3 的定位结果，分别报告机制命中和计时判定。

工具输出只证明给定输入下的行为。Memcheck 的 0 命中不覆盖算术泄漏与缓存泄漏。Binsec 的 bounded secure 只覆盖给定入口、长度与秘密模型。Microwalk 的互信息只覆盖实际执行的用例。dudect 与 RTLF 只回答分布是否可区分。把这些边界写进表，是为了让“未运行”与“未检出”在版面上就分开，而不是靠读者猜。

\section{90 格排查全表}

本附录给出 3 实现 × 5 步骤 × 6 类 = 90 格的逐格判定。来源列的数字含义：1 敏感分支/提前返回、2 可变时延指令/比较、3 数据相关循环、4 数据相关访存、5 失败路径、6 异常处理。判定只取“存在 / 不存在 / 不适用 / 未检查”。全表合计：存在 5、不存在 75、不适用 10、未检查 0。人工夹具 M3、M4 不计入本表。“存在”格中的 M5 一格携带公开信息，按有界例外登记。

\begin{center}
{\scriptsize
\renewcommand{\arraystretch}{0.95}
\begin{tabularx}{\textwidth}{|c|c|>{\raggedright\arraybackslash}p{2.15cm}|c|>{\raggedright\arraybackslash}p{3.2cm}|>{\raggedright\arraybackslash}p{3.0cm}|>{\centering\arraybackslash}X|}
\hline
\apptabhead{实现} & \apptabhead{步骤} & \apptabhead{来源} & \apptabhead{判定} & \apptabhead{函数} & \apptabhead{指令或行号} & \apptabhead{证据} \\
\hline
M1 & 解密 & 1 敏感分支/提前返回 & 不存在 & indcpa\_dec / poly\_tomsg & 常数回边 & A-17 \\
M1 & 解密 & 2 可变时延指令/比较 & 不存在 & poly\_tomsg & 乘加、移位，无 \texttt{div} & A-17 \\
M1 & 解密 & 3 数据相关循环 & 不存在 & poly\_tomsg & 32 × 8 常数上界 & A-17 \\
M1 & 解密 & 4 数据相关访存 & 不存在 & poly\_tomsg & 固定顺序写 & A-17 \\
M1 & 解密 & 5 失败路径 & 不存在 & crypto\_kem\_dec & 本步不判定失败 & A-17 \\
M1 & 解密 & 6 异常处理 & 不适用 & — & 无异常机制 & A-17 \\
\hline
M1 & 哈希 & 1 敏感分支/提前返回 & 不存在 & hash\_g / hash\_h / rkprf & 定长调用 & A-17 \\
M1 & 哈希 & 2 可变时延指令/比较 & 不存在 & FIPS202 & 无比较 & A-17 \\
M1 & 哈希 & 3 数据相关循环 & 不存在 & FIPS202 & 24 轮固定 & A-17 \\
M1 & 哈希 & 4 数据相关访存 & 不存在 & FIPS202 & 常量下标 & A-17 \\
M1 & 哈希 & 5 失败路径 & 不存在 & FIPS202 & 不返回状态 & A-17 \\
M1 & 哈希 & 6 异常处理 & 不存在 & FIPS202 & 无 abort/断言 & A-17 \\
\hline
M1 & 重加密 & 1 敏感分支/提前返回 & 存在 & poly\_frommsg & \texttt{0x16e bt} → \texttt{0x172 jae} & A-17、A-22 \\
M1 & 重加密 & 2 可变时延指令/比较 & 不存在 & poly\_frommsg / 压缩 & 无比较动作 & A-17 \\
M1 & 重加密 & 3 数据相关循环 & 不存在 & poly\_frommsg & 32 × 8 常数 & A-17 \\
M1 & 重加密 & 4 数据相关访存 & 不存在 & poly\_frommsg & 固定顺序写 & A-17 \\
M1 & 重加密 & 5 失败路径 & 不存在 & indcpa\_enc & 不返回错误 & A-17 \\
M1 & 重加密 & 6 异常处理 & 不存在 & indcpa\_enc & 无异常机制 & A-17 \\
\hline
M1 & 比较 & 1 敏感分支/提前返回 & 不存在 & verify & 遍历 1088 B & A-17 \\
M1 & 比较 & 2 可变时延指令/比较 & 不存在 & verify & 累积异或，不短路 & A-17 \\
M1 & 比较 & 3 数据相关循环 & 不存在 & verify & 长度常数 1088 & A-17 \\
M1 & 比较 & 4 数据相关访存 & 不存在 & verify & 顺序下标 & A-17 \\
M1 & 比较 & 5 失败路径 & 不存在 & verify & 结果只作掩码 & A-17 \\
M1 & 比较 & 6 异常处理 & 不适用 & — & 不可失败 & A-17 \\
\hline
M1 & 隐式拒绝 & 1 敏感分支/提前返回 & 不存在 & dec / cmov & 无条件走完 & A-17 \\
M1 & 隐式拒绝 & 2 可变时延指令/比较 & 不存在 & cmov & \texttt{0x2b26 sete} → \texttt{0x2b37 call} & A-17 \\
M1 & 隐式拒绝 & 3 数据相关循环 & 不存在 & rkprf / cmov & 定长 32 B & A-17 \\
M1 & 隐式拒绝 & 4 数据相关访存 & 不存在 & cmov & 固定偏移、顺序下标 & A-17 \\
M1 & 隐式拒绝 & 5 失败路径 & 不存在 & dec & 单点返回，零跳转 & A-17 \\
M1 & 隐式拒绝 & 6 异常处理 & 不适用 & — & 无异常机制 & A-17 \\
\hline
\multicolumn{7}{|l|}{M1 小计：存在 1、不存在 26、不适用 3、未检查 0（30 格）} \\
\hline
\end{tabularx}
}
\apptabnote{表 F1　M1 的 30 格排查。}
\end{center}

\begin{center}
{\scriptsize
\renewcommand{\arraystretch}{0.95}
\begin{tabularx}{\textwidth}{|c|c|>{\raggedright\arraybackslash}p{2.15cm}|c|>{\raggedright\arraybackslash}p{3.2cm}|>{\raggedright\arraybackslash}p{3.0cm}|>{\centering\arraybackslash}X|}
\hline
\apptabhead{实现} & \apptabhead{步骤} & \apptabhead{来源} & \apptabhead{判定} & \apptabhead{函数} & \apptabhead{指令或行号} & \apptabhead{证据} \\
\hline
M2 & 解密 & 1 敏感分支/提前返回 & 不存在 & dec / indcpa\_dec & 常数回边 & A-18 \\
M2 & 解密 & 2 可变时延指令/比较 & 存在 & poly\_tomsg & \texttt{0x36bf div \%r9d} & A-18、A-22 \\
M2 & 解密 & 3 数据相关循环 & 不存在 & poly\_tomsg & 32 × 8 常数 & A-18 \\
M2 & 解密 & 4 数据相关访存 & 不存在 & poly\_tomsg & 顺序读写 & A-18 \\
M2 & 解密 & 5 失败路径 & 不存在 & crypto\_kem\_dec & 本步不判定失败 & A-18 \\
M2 & 解密 & 6 异常处理 & 不适用 & — & RNG 不在本步 & A-18 \\
\hline
M2 & 哈希 & 1 敏感分支/提前返回 & 不存在 & sha3 / shake & 定长调用 & A-18 \\
M2 & 哈希 & 2 可变时延指令/比较 & 不存在 & FIPS202 & 无比较 & A-18 \\
M2 & 哈希 & 3 数据相关循环 & 不存在 & FIPS202 & 24 轮固定 & A-18 \\
M2 & 哈希 & 4 数据相关访存 & 不存在 & FIPS202 & 常量下标 & A-18 \\
M2 & 哈希 & 5 失败路径 & 不存在 & FIPS202 & 不返回状态 & A-18 \\
M2 & 哈希 & 6 异常处理 & 不存在 & FIPS202 & 无 abort/断言 & A-18 \\
\hline
M2 & 重加密 & 1 敏感分支/提前返回 & 存在 & poly\_frommsg & \texttt{0x3668 bt} → \texttt{0x366c jae} & A-23、A-22 \\
M2 & 重加密 & 2 可变时延指令/比较 & 存在 & poly\_compress / polyvec\_compress & \texttt{0x3523 div}、\texttt{0x38c7 div} & A-18、A-22 \\
M2 & 重加密 & 3 数据相关循环 & 不存在 & 压缩函数 & 常数上界 & A-18 \\
M2 & 重加密 & 4 数据相关访存 & 不存在 & 压缩函数 & 定长位移拼装 & A-18 \\
M2 & 重加密 & 5 失败路径 & 不存在 & indcpa\_enc & 不返回错误 & A-18 \\
M2 & 重加密 & 6 异常处理 & 不存在 & indcpa\_enc & 无异常机制 & A-18 \\
\hline
M2 & 比较 & 1 敏感分支/提前返回 & 不存在 & verify & 遍历 1088 B & A-18 \\
M2 & 比较 & 2 可变时延指令/比较 & 不存在 & verify & 累积异或 & A-18 \\
M2 & 比较 & 3 数据相关循环 & 不存在 & verify & 常数 1088 & A-18 \\
M2 & 比较 & 4 数据相关访存 & 不存在 & verify & 顺序下标 & A-18 \\
M2 & 比较 & 5 失败路径 & 不存在 & verify / cmov & 结果只作掩码 & A-18 \\
M2 & 比较 & 6 异常处理 & 不适用 & — & 不可失败 & A-18 \\
\hline
M2 & 隐式拒绝 & 1 敏感分支/提前返回 & 不存在 & dec / cmov & 无条件走完 & A-18 \\
M2 & 隐式拒绝 & 2 可变时延指令/比较 & 不存在 & cmov & 值屏障在位 & A-18 \\
M2 & 隐式拒绝 & 3 数据相关循环 & 不存在 & cmov / KDF & 32 B 与定长 & A-18 \\
M2 & 隐式拒绝 & 4 数据相关访存 & 不存在 & cmov & 固定偏移 & A-18 \\
M2 & 隐式拒绝 & 5 失败路径 & 不存在 & dec & 同一路径，返回 0 & A-18 \\
M2 & 隐式拒绝 & 6 异常处理 & 不适用 & — & 无异常机制 & A-18 \\
\hline
\multicolumn{7}{|l|}{M2 小计：存在 3、不存在 24、不适用 3、未检查 0（30 格）} \\
\hline
\end{tabularx}
}
\apptabnote{表 F2　M2 的 30 格排查。重加密的“存在”格在未修复版与已修复版两侧均成立。}
\end{center}

\begin{center}
{\scriptsize
\renewcommand{\arraystretch}{0.95}
\begin{tabularx}{\textwidth}{|c|c|>{\raggedright\arraybackslash}p{2.15cm}|c|>{\raggedright\arraybackslash}p{3.2cm}|>{\raggedright\arraybackslash}p{3.0cm}|>{\centering\arraybackslash}X|}
\hline
\apptabhead{实现} & \apptabhead{步骤} & \apptabhead{来源} & \apptabhead{判定} & \apptabhead{函数} & \apptabhead{指令或行号} & \apptabhead{证据} \\
\hline
M5 & 解密 & 1 敏感分支/提前返回 & 不存在 & indcpa\_dec & 常数回边 & A-21 \\
M5 & 解密 & 2 可变时延指令/比较 & 不存在 & poly\_tomsg & 乘加、移位 & A-21 \\
M5 & 解密 & 3 数据相关循环 & 不存在 & indcpa\_dec & 常数上界 & A-21 \\
M5 & 解密 & 4 数据相关访存 & 不存在 & indcpa\_dec & 计数器下标 & A-21 \\
M5 & 解密 & 5 失败路径 & 不存在 & indcpa\_dec & 仅分配失败清理 & A-21 \\
M5 & 解密 & 6 异常处理 & 不适用 & — & 无异常机制 & A-21 \\
\hline
M5 & 哈希 & 1 敏感分支/提前返回 & 不存在 & hash\_g / hash\_h / hash\_j & 定长调用 & A-21 \\
M5 & 哈希 & 2 可变时延指令/比较 & 不存在 & FIPS202 & 长度驱动 & A-21 \\
M5 & 哈希 & 3 数据相关循环 & 不存在 & FIPS202 & 固定轮数 & A-21 \\
M5 & 哈希 & 4 数据相关访存 & 不存在 & FIPS202 & 常量下标 & A-21 \\
M5 & 哈希 & 5 失败路径 & 不存在 & FIPS202 & 不返回状态 & A-21 \\
M5 & 哈希 & 6 异常处理 & 不存在 & FIPS202 & 无 abort/断言 & A-21 \\
\hline
M5 & 重加密 & 1 敏感分支/提前返回 & 不存在 & poly\_frommsg & 唯一跳转为循环回边 & A-21 \\
M5 & 重加密 & 2 可变时延指令/比较 & 不存在 & 压缩函数 & 乘加、位移 & A-21 \\
M5 & 重加密 & 3 数据相关循环 & 不存在 & poly\_frommsg & 32 次定长 & A-21 \\
M5 & 重加密 & 4 数据相关访存 & 不存在 & poly\_frommsg & 固定顺序 & A-21 \\
M5 & 重加密 & 5 失败路径 & 不存在 & indcpa\_enc & 仅分配失败清理 & A-21 \\
M5 & 重加密 & 6 异常处理 & 不适用 & — & 无异常机制 & A-21 \\
\hline
M5 & 比较 & 1 敏感分支/提前返回 & 不存在 & ct\_memcmp & 循环体内无跳转 & A-21 \\
M5 & 比较 & 2 可变时延指令/比较 & 不存在 & ct\_memcmp & 无分支归约 & A-21 \\
M5 & 比较 & 3 数据相关循环 & 不存在 & ct\_memcmp & 常数 1088 & A-21 \\
M5 & 比较 & 4 数据相关访存 & 不存在 & ct\_memcmp & 顺序 64 B/轮 & A-21 \\
M5 & 比较 & 5 失败路径 & 不存在 & ct\_memcmp & 返回不构造分支 & A-21 \\
M5 & 比较 & 6 异常处理 & 不适用 & — & 不可失败 & A-21 \\
\hline
M5 & 隐式拒绝 & 1 敏感分支/提前返回 & 存在（公开信息） & check\_sk & \texttt{a404 call} → \texttt{a40e jne ad87} & A-21、A-22 \\
M5 & 隐式拒绝 & 2 可变时延指令/比较 & 不存在 & ct\_cmov\_zero & \texttt{ad60} 至 \texttt{ad85} 掩码 & A-21 \\
M5 & 隐式拒绝 & 3 数据相关循环 & 不存在 & KDF / cmov & 固定轮数与 32 B & A-21 \\
M5 & 隐式拒绝 & 4 数据相关访存 & 不存在 & KDF / cmov & 固定偏移 & A-21 \\
M5 & 隐式拒绝 & 5 失败路径 & 不存在 & dec & 成功与失败同路径 & A-21 \\
M5 & 隐式拒绝 & 6 异常处理 & 不适用 & — & 无异常机制 & A-21 \\
\hline
\multicolumn{7}{|l|}{M5 小计：存在 1、不存在 25、不适用 4、未检查 0（30 格）} \\
\hline
\end{tabularx}
}
\apptabnote{表 F3　M5 的 30 格排查。唯一“存在”格判据是公开数据的哈希，按有界例外登记。}
\end{center}

三张表合计 90 格：存在 5、不存在 75、不适用 10、未检查 0。五格“存在”分别是 M1 重加密的敏感条件分支、M2 解密的可变时延指令（除法）、M2 重加密的可变时延指令（除法）、M2 重加密的敏感条件分支（未修复版与已修复版两侧），以及 M5 隐式拒绝的公开 sk 检查。十格“不适用”全部是异常处理类与该步骤不可能失败的组合。

90 格按“存在 / 不存在 / 不适用 / 未检查”分别登记。M5 的“存在”格是公开 sk 完整性检查，判据是公开数据的哈希，按有界例外登记。排查同时发现 M2 官方修复产物中仍保留与 M1 同源的消息编码分支，这是全流程方法对补丁产物的超额发现。

五个“存在”格分别落在三条机制上：M1 的消息编码分支、M2 的三处秘密除法与消息编码分支、M5 的公开 sk 检查。前四条与秘密相关，第五条只携带公开信息。把五格都写成“漏洞”会混淆可利用性，把第五条删掉又会掩盖一个有意的提前退出。表 F1 至表 F3 逐格保留函数与地址，是为了让读者能回到产物核对，而不是只看一个总数。

十格“不适用”的分布也要能对上：M1 三格、M2 三格、M5 四格，全部是异常处理类与该步骤不可能失败的组合。它们不是“没查”，也不是“查了没有”，而是该步骤没有这类机制可查。三种判定各自计数：“不存在”只用于已查无发现，“不适用”只用于无此类机制，“未检查”单列。这样“不适用”既不会并入阴性，也不会并入缺口。M3 与 M4 是人工校准，不计入 90 格，它们在正文与附录 E 单列，避免与真实实现混算。

\section{复现入口与命令}

本附录的命令均在 x86\_64 Linux 用户态实跑。不写 sysfs、不需要 sudo。实跑环境为 WSL2 虚拟化 Linux（clang 14、CMake 3.31、OpenSSL 3），与锁频采数所用的裸金属不是同一台机器。命令的通过数用于验证软件流程，不产生计时结论。完整锁频采集需要 x86\_64 裸金属与频率写权限，不在本附录的实跑范围内。

\begin{center}
{\footnotesize
\renewcommand{\arraystretch}{1.1}
\begin{tabularx}{\textwidth}{|>{\raggedright\arraybackslash}p{2.55cm}|>{\raggedright\arraybackslash}p{5.15cm}|>{\raggedright\arraybackslash}p{2.55cm}|>{\raggedright\arraybackslash}X|}
\hline
\apptabhead{目标} & \apptabhead{命令} & \apptabhead{预期输出} & \apptabhead{实测结果（权限 / 机器）} \\
\hline
普通构建与测试 & \texttt{make test} & 普通 CTest 全通过 & 57/57 通过，0 失败（普通用户。WSL2，clang 14） \\
\hline
消毒器构建与测试 & \texttt{make sanitize} & ASan/UBSan CTest 全通过 & 57/57 通过，0 失败（普通用户。WSL2，clang 14。ASan 关闭泄漏检测并关闭地址随机化） \\
\hline
平台配置 & \texttt{cmake -S . -B build/platform -DBASELINE\_ENABLE\_PLATFORM=ON -DA113\_ARCHIVED\_CASE\_ONLY=OFF} & 配置成功，生成构建树 & 配置成功（普通用户，WSL2） \\
\hline
平台构建 & \texttt{cmake --build build/platform} & 构建成功 & 构建成功（普通用户，WSL2） \\
\hline
平台测试 & \texttt{ctest --test-dir build/platform --output-on-failure} & 平台 CTest 全通过 & 87/87 通过，0 失败（普通用户。vendor 未启用。quick 档启用后 88/88） \\
\hline
一键 quick & \texttt{sh scripts/reproduce\_all.sh quick --fetch} & 四配置构建与 CTest 全通过，统计演练完成 & 263/263 通过，0 失败（88 + 57 + 57 + 61。普通用户，WSL2） \\
\hline
vendor 取件 & \texttt{sh tools/fetch\_all\_vendors.sh} & 四棵 vendor 树就位并逐项 SHA-256 校验 & 4/4 就位，摘要校验通过（普通用户，WSL2） \\
\hline
full 档演练 & \texttt{sh scripts/reproduce\_all.sh full --modules M2,M3,M4,M5 --control-root <控制根>} & 只打印采集命令，不执行 & 打印成功，退出码 0（普通用户，WSL2） \\
\hline
\end{tabularx}
}
\apptabnote{表 G1　复现入口与命令。命令实跑通过。完整锁频采集需要 x86\_64 裸金属与频率写权限。quick 档的 263 项是 platform 88 + ordinary 57 + sanitize 57 + historical-fixed 61，失败 0。}
\end{center}

quick 档只做配置、构建、CTest 与小样本统计演练：统计演练每类 200 个样本、1 次重复，用于检查链路是否连通。它不锁频、不写 sysfs、不需要 sudo、不绑核，验证软件流程与测试入口，不产生任何模块的正式统计结论。full 档按冻结计划生成采集命令，默认只打印，加 \texttt{--execute} 才真正采集，且必须给出控制根。计时结论来自锁频采集，仍需裸金属与频率写权限。干净 clone 缺 vendor 与 third\_party，M1 与 M2 的 vendor 树缺省不存在，先跑取件脚本。取件按固定摘要校验，不覆盖已有且摘要一致的树。

实跑环境的编译器与采数环境不同：本附录的命令在 WSL2 的 clang 14 与 CMake 3.22 下执行，锁频数据来自另一台 x86\_64 裸金属的 clang 18.1.3。两边只共享源码与测试入口，不共享时间结论。把这一点写在表 G1 的表注里，是为了避免读者把 263 项通过当成 263 次计时实验。

命令失败时的排查顺序也写进复现说明：先看编译器与 CMake 版本，再看 vendor 是否就位，然后看 sanitizer 的运行库。ASan 在部分虚拟化环境需要关闭地址随机化或关闭泄漏检测才能跑完。这类环境差异只影响测试入口，不影响源码。复现者若在另一台机器上得到不同的通过数，应先核对配置项与编译器，而不是先怀疑实验结论。



\section{术语与缩略语}

表 H1 给出全篇统一术语。术语首次出现时列出中文说明与英文全称，之后统一使用定稿词。

\begin{center}
{\footnotesize
\renewcommand{\arraystretch}{1.0}
\begin{tabularx}{\textwidth}{|>{\raggedright\arraybackslash}p{1.85cm}|>{\raggedright\arraybackslash}p{4.0cm}|>{\raggedright\arraybackslash}X|>{\raggedright\arraybackslash}p{1.5cm}|}
\hline
\apptabhead{术语} & \apptabhead{英文 / 全称} & \apptabhead{含义} & \apptabhead{首见} \\
\hline
ML-KEM & Module-Lattice-Based Key-Encapsulation Mechanism & FIPS 203 标准化的密钥封装机制。本作品的算法标准为 ML-KEM-768。 & 1.1 \\
\hline
Kyber & CRYSTALS-Kyber & ML-KEM 的前身。本作品的 M2 是 pq-crystals 历史快照 Kyber768。 & 1.1 \\
\hline
FO & Fujisaki-Okamoto & 把公钥加密转换为 KEM 的变换。解封装中的重加密与比较由它规定。 & 1.1 \\
\hline
KEM & Key-Encapsulation Mechanism & 密钥封装机制。 & 1.1 \\
\hline
PQC & Post-Quantum Cryptography & 后量子密码。 & 1.1 \\
\hline
MDE & Minimum Detectable Effect & 最小可检测效应。给定 $\alpha$、功效、样本量与方差下可分辨的最小均值差。 & 2.4 \\
\hline
Holm & Holm-Bonferroni & 逐步下降的多重比较校正。本作品四个家族各取全族 $\alpha=0.05$。 & 2.4 \\
\hline
RTLF & R-Time-Leak-Finder & 基于分位数与 bootstrap 的族外分布检验。$\alpha=0.01$。 & 2.3 \\
\hline
dudect & Dude, is my code constant time? & 基于 Welch $t$ 的在线泄漏检测工具。最大 $t$ 不转成 $p$ 值。 & 2.3 \\
\hline
Binsec & Binsec/Rel & 二进制级关系符号执行与有界常数时间检查工具。 & 2.3 \\
\hline
Microwalk & MicroWalk & 基于动态二进制插桩的访存互信息分析工具。 & 2.3 \\
\hline
Valgrind & Valgrind & 动态分析框架。本作品使用 Memcheck 与 varlat 补丁。 & 2.3 \\
\hline
varlat & variable-latency patch & Valgrind 的可变时延指令检查补丁，用于除法操作数含秘密的情形。 & 2.3 \\
\hline
aarch64 & ARM 64-bit architecture & ARM 64 位架构。 & 4.11 \\
\hline
x86\_64 & x86-64 & 本作品的主测量架构。 & 1.1 \\
\hline
ipd & initial public draft & FIPS 203 初始公开草案。M1 的对象口径。 & 1.3 \\
\hline
\end{tabularx}
}
\apptabnote{表 H1　术语与缩略语。术语全篇统一。第一次出现给全称，之后只用定稿词。}
\end{center}

\fussy
